% ---------------------------------------------------------------------------
% Chapter 1: The Ancestral Home: Protons, PDFs and the Hard Scattering
% Source files quoted in this chapter live in example/ME/ and are read at
% build time (minted), so the book always shows the current version.
% ---------------------------------------------------------------------------
\chapter{The Ancestral Home: Protons, PDFs and the Hard Scattering}
\chaptermark{The Ancestral Home}
\label{ch:ancestral_home}

\begin{journeybox}
Every family has an ancestral home. Ours is the proton, and the mother of our
jet, a bottom quark, is born there in a hard scattering. Before she takes a
single step we must know how much energy she was given, because everything that
follows is measured against it.
\end{journeybox}

\physicsline{factorisation, PDFs, the leading-order matrix element for
$q\bar q\to\bbbar$, scale choices, \PYTHIA's \code{HardQCD:qqbar2bbbar} hard
process inside \CMSSW, the hard-process record of our event}

\section{What this chapter computes, and why we do it three times}
\label{sec:ancestral_home:what_we_compute}

The whole chapter is about one number: the cross section for producing a
$\bbbar$ pair with $\pthat\ge 30\GeV$ in proton--proton collisions at
$\sqrt{s}=13.6\TeV$ through quark--antiquark annihilation. It comes out at
about $18.5\nb$. We compute it three times, each time with less of our own
work and more trust in somebody else's code:
\begin{enumerate}[nosep]
  \item \textbf{Theory.} From the factorisation formula, the parton
    distribution functions and a one-line matrix element we write down a
    three-dimensional integral and evaluate it
    (\crefrange{sec:ancestral_home:two_protons_meet}{sec:ancestral_home:from_formula_to_number}).
  \item \textbf{Standalone code.} A single Python file of about 190 lines,
    needing only NumPy and SciPy, does exactly that integral. The reader runs it
    and gets the same number in about one second
    (\cref{sec:ancestral_home:standalone_code}).
  \item \textbf{\CMSSW.} The official CMS generator configuration runs the same
    process inside \PYTHIA. Its log prints the cross section it sampled
    (\cref{sec:ancestral_home:cmssw}).
\end{enumerate}
\Cref{tab:ancestral_home:three_numbers} anticipates the result. The three
agree within the statistical precision of the Monte Carlo runs, and that
agreement is the point of the exercise: after this chapter the reader knows
what every knob in the CMS configuration does to the number, because they have
turned each knob themselves.

\begin{table}[htbp]\centering\small
\caption[The $q\bar q\to\bbbar$ cross section obtained three ways]{The $q\bar q\to\bbbar$ cross section at $\sqrt{s}=13.6\TeV$ with
$\pthat\ge30\GeV$, $\mb=4.8\GeV$, NNPDF3.1 NNLO and
$\muR^2=\muF^2=\pthat^2+\mb^2$, obtained three ways. The scripts are
\file{standalone_qqbar_bbbar_xsec.py} and \file{qqbar_bbbar_xsec.py} in
\file{example/ME/Standalone/}; the \CMSSW number is from
\file{example/ME/CMSSW/step1_GEN_cfg.log}. Quoted errors are statistical or
numerical only; the theoretical uncertainty is discussed in
\cref{sec:ancestral_home:scales}.}
\label{tab:ancestral_home:three_numbers}
\begin{tabular}{llr}\toprule
Method & Details & $\sigma$ [nb] \\ \midrule
Standalone script, one pass & LO integral, $2^{16}$ Sobol points & $18.5064$ \\
Standalone script, 8 repeats & LO integral, scrambled seeds & $18.50653 \pm 0.00005$ \\
Full script (LHAPDF, \PYTHIA $\alphas$) & LO integral, 8 repeats & $18.50662 \pm 0.00005$ \\
\PYTHIA~8.315 standalone & Monte Carlo, large sample & $18.44 \pm 0.10$ \\
\PYTHIA~8.311 in \code{CMSSW\_15\_0\_5} & Monte Carlo, 1000 events & $18.33 \pm 0.30$ \\
\bottomrule
\end{tabular}
\end{table}

A word on the tools. The preface promises that \MG generates Event~A. For
this first stage we use \PYTHIA's own $2\to2$ matrix element instead,
because it is the process CMS actually runs for the fragment in
\cref{sec:ancestral_home:cmssw}, and because its leading-order matrix element is
a formula short enough to check by hand. \MG at leading order uses the same
Born matrix element; the correspondence between the two sets of input cards is
given in \cref{sec:ancestral_home:madgraph}.

\section{Two protons meet: the factorisation theorem in one page}
\label{sec:ancestral_home:two_protons_meet}

A proton is not a particle one can scatter off in perturbation theory. The
collinear factorisation theorem separates the problem into a long-distance
part that is measured once and for all, the parton distribution functions
(PDFs), and a short-distance part that is calculated, the partonic cross
section. For any inclusive hard process $pp\to\bbbar+X$,
\begin{equation}
\begin{split}
  \sigma(pp\to\bbbar X) &= \sum_{i,j}\int_0^1\!\mathrm{d}x_1\int_0^1\!\mathrm{d}x_2\;
  f_i(x_1,\muF^2)\, f_j(x_2,\muF^2)\;
  \hat\sigma_{ij\to\bbbar}\!\left(\shat;\, \muR^2,\muF^2\right) \\
  &\quad + \mathcal{O}\!\left(\frac{\Lambda^2_{\mathrm{QCD}}}{\pthat^2}\right),
  \qquad \shat = x_1x_2 s .
\end{split}
  \label{eq:ancestral_home:factorisation}
\end{equation}
Here $f_i(x,\muF^2)\,\mathrm{d}x$ is the number of partons of flavour $i$
carrying a fraction between $x$ and $x+\mathrm{d}x$ of the proton momentum,
resolved at the factorisation scale $\muF$; $\hat\sigma_{ij}$ is the cross
section for the two partons, computed in perturbation theory with the strong
coupling evaluated at the renormalisation scale $\muR$; and $\shat$ is the
squared partonic centre-of-mass energy. With $\sqrt{s}=13.6\TeV$ each beam
carries $6800\GeV$. The two protons are identical, so the sum over $(i,j)$
counts both $q$ from beam~1 with $\bar q$ from beam~2 and the reverse; both
orientations must be included.

\paragraph{Better variables.} For a $2\to2$ process the natural variables
are the pair's squared-mass fraction and rapidity,
\begin{equation}
  \tau \equiv x_1x_2 = \frac{\shat}{s},\qquad
  Y \equiv \tfrac12\ln\frac{x_1}{x_2},\qquad
  x_{1,2} = \sqrt{\tau}\,e^{\pm Y}.
  \label{eq:ancestral_home:tauY}
\end{equation}
The Jacobian is $|\partial(x_1,x_2)/\partial(\ln\tau,Y)| = \tau$, hence
$\mathrm{d}x_1\,\mathrm{d}x_2 = \tau\,\mathrm{d}\ln\tau\,\mathrm{d}Y$. This is
more than a convenience. PDF libraries return $x f(x,Q^2)$, not $f$, and
\begin{equation}
  f_i(x_1)f_j(x_2)\,\mathrm{d}x_1\mathrm{d}x_2
  = \frac{[x_1f_i(x_1)]\,[x_2f_j(x_2)]}{\tau}\,\tau\,\mathrm{d}\ln\tau\,\mathrm{d}Y
  = [x_1f_i(x_1)]\,[x_2f_j(x_2)]\;\mathrm{d}\ln\tau\,\mathrm{d}Y ,
  \label{eq:ancestral_home:xf}
\end{equation}
so the factors of $\tau$ cancel exactly and the code never divides by $x$.
The limits are $\tau_{\min}\le\tau\le1$ with $\tau_{\min}$ fixed by the
kinematic cut (\cref{sec:ancestral_home:from_formula_to_number}), and
$|Y|\le -\tfrac12\ln\tau$ so that both $x_1$ and $x_2$ stay below one.

\section{Parton distribution functions: where the mother comes from}
\label{sec:ancestral_home:pdfs}

The PDFs cannot be computed from first principles; they are fitted to data
and evolved in $Q^2$ with the DGLAP equations. CMS Run~3 simulation uses the
NNPDF3.1 NNLO set with $\alphas(m_Z)=0.118$~\cite{NNPDF:2017mvq}, member~0
(the central replica), accessed through LHAPDF~\cite{Buckley:2014ana} under the
name \code{NNPDF31\_nnlo\_as\_0118} and the numeric ID~303600. The set is
distributed as a table of $x f(x,Q^2)$ on a grid of knots in $x$ and $Q$;
LHAPDF interpolates between the knots with cubic splines in $(\ln x,\ln Q^2)$.

Three conventions of the interface matter for reproducing \PYTHIA's numbers,
and the reader would not guess any of them from the physics:
\begin{itemize}[nosep]
  \item \textbf{Freezing at the grid edges.} Outside the tabulated range in
    $x$ or $Q^2$ the value at the nearest edge is used rather than an
    extrapolation.
  \item \textbf{Clipping.} Fitted PDFs can go slightly negative at large $x$
    and for heavy flavours; \PYTHIA's proton wrapper sets such values to zero.
    In the default integration $1.17$ million of the individual PDF evaluations
    were clipped, yet the effect on the cross section is $0.05\%$
    (\cref{tab:ancestral_home:variations}), because the clipped values are tiny.
  \item \textbf{Symmetric heavy sea.} The charm and bottom antiquark
    densities are taken equal to the quark densities, $\bar c=c$ and
    $\bar b=b$, as in the wrapper.
\end{itemize}

\paragraph{Our event.} The hard-process record of the event followed through
this book (\cref{sec:ancestral_home:our_jet}) was produced by an antiquark
from beam~A and a quark from beam~B, and \PYTHIA prints the PDF values it used
(\cref{tab:ancestral_home:pdf_values}). The two partons could hardly be more
different. The $\bar u$ carries $x_1=2.75\times10^{-4}$ of its proton: it is a
sea antiquark, drawn from the steeply rising small-$x$ region where $xf$ is of
order one. The $u$ carries $x_2=0.284$: a valence quark from the region where
$xf\approx0.4$. The asymmetry is the first fact about our mother's birth.
The pair is boosted along the beam with rapidity $Y=\tfrac12\ln(x_1/x_2)=-3.47$,
so both $b$ quarks emerge far backward, and the mother we follow is born
pointing at the very edge of the CMS tracker.

\begin{table}[htbp]\centering
\caption[PDF values and coupling at our event]{The PDF values and coupling at our event, as printed by \PYTHIA in
\file{example/ME/CMSSW/step1_GEN_cfg.log}, compared with an independent
evaluation of the same LHAPDF set. The scale is
$Q^2=\pthat^2+\mb^2=1178\GeV^2$. The small differences come from the rounded
$x$ and $Q^2$ in the printout.}
\label{tab:ancestral_home:pdf_values}
\begin{tabular}{llrr}\toprule
Parton & $x$ & $xf$ (\PYTHIA printout) & $xf$ (LHAPDF, this book) \\ \midrule
$\bar u$ from beam A & $2.752\times10^{-4}$ & $1.922$ & $1.92184$ \\
$u$ from beam B & $0.2841$ & $0.4075$ & $0.40742$ \\ \midrule
$\alphas(Q^2)$ & & $0.1385$ & $0.138534$ \\
\bottomrule
\end{tabular}
\end{table}

\section{The matrix element for \texorpdfstring{$\bbbar$}{bb} production}
\label{sec:ancestral_home:matrix_element}

\subsection{The \texorpdfstring{$q\bar q\to\bbbar$}{qq->bb} channel at leading order}

The Born process is a single $s$-channel gluon exchange. Averaging over the
spins and colours of the incoming massless quark pair and summing over those
of the massive outgoing pair, the differential cross section in the partonic
centre-of-mass frame is~\cite{ParticleDataGroup:2024cfk}
\begin{equation}
  \frac{\mathrm{d}\hat\sigma}{\mathrm{d}\cos\theta}
  = \frac{\pi\alphas^2(\muR^2)\,\beta}{9\,\shat}
    \left[\,1+\cos^2\theta+\rho\,(1-\cos^2\theta)\,\right],
  \qquad
  \rho\equiv\frac{4\mb^2}{\shat},\quad \beta\equiv\sqrt{1-\rho},
  \label{eq:ancestral_home:dsigma}
\end{equation}
where $\theta$ is the angle between the incoming quark and the outgoing $b$,
$\beta$ is the $b$ velocity and $\rho$ vanishes far above threshold. The
$1/9$ is the colour factor: the average $1/3\cdot1/3$ over incoming colours
times the sum $\mathrm{Tr}(T^aT^b)\mathrm{Tr}(T^aT^b)=2$ over the gluon
exchange, combined with the spin average. Integrating over $\cos\theta$ from
$-1$ to $1$ gives the total partonic cross section
\begin{equation}
  \hat\sigma(\shat) = \frac{4\pi\alphas^2\,\beta\,(2+\rho)}{27\,\shat},
  \label{eq:ancestral_home:sigmahat}
\end{equation}
which reduces to the familiar $8\pi\alphas^2/(27\shat)$ for massless quarks.
This is exactly the function \code{sigmaKin} of \PYTHIA's class
\code{Sigma2qqbar2QQbar}, process code~124. At our event,
$\shat=14\,456\GeV^2$ and $\alphas=0.1385$ give $\hat\sigma=481\pb$ and
$\mathrm{d}\hat\sigma/\mathrm{d}\cos\theta=302\pb$ at the actual
$\cos\hat\theta=-0.824$; the reader is asked to verify both in
\cref{ch:ancestral_home} Exercise~2. The conversion from natural units is
$(\hbar c)^2 = 0.389\,379\,372\GeV^2\,\mathrm{mb} = 3.893\,793\,72\times10^{8}\GeV^2\,\mathrm{pb}$.

\paragraph{Kinematics of the two quarks.} In the partonic frame each $b$ has
momentum $p^{*}=\sqrt{\shat/4-\mb^2}$, and
\begin{equation}
  \pthat^2 = p^{*2}\,(1-\cos^2\theta),\qquad
  y^{*} = \operatorname{artanh}(\beta\cos\theta),\qquad
  y_{b,\bar b} = Y \pm y^{*},\qquad
  \eta = \operatorname{arsinh}\!\left(\frac{m_{\mathrm{T}}}{\pt}\sinh y\right),
  \label{eq:ancestral_home:kinematics}
\end{equation}
with $m_{\mathrm{T}}^2=\pt^2+\mb^2$. At Born level the two quarks have equal
and opposite transverse momentum, so ``$\pthat$'' is the \pt of either. A cut
$\pthat\ge\pthat^{\min}$ is therefore a cut on the angle,
$|\cos\theta|\le c_{\max}=\sqrt{1-(\pthat^{\min}/p^{*})^2}$, and it can only
be satisfied when $p^{*}>\pthat^{\min}$, that is when
$\shat\ge 4(\mb^2+\pthat^{\min\,2})$. This sets $\tau_{\min}$ in
\cref{eq:ancestral_home:tauY}: for $\pthat^{\min}=30\GeV$ and $\mb=4.8\GeV$,
$\tau_{\min}=4\times(4.8^2+30^2)/13600^2 = 2.0\times10^{-5}$. Note that
rapidity $y$ and pseudorapidity $\eta$ differ for a massive quark; the
difference is $0.01$ at $\pt=34\GeV$ and would be far larger near threshold.

\subsection{The \texorpdfstring{$gg$}{gg} channel, the incoming \texorpdfstring{$b$}{b}, and four versus five flavours}
\label{sec:ancestral_home:gg_and_flavours}

At LHC energies the gluon luminosity is so large that $gg\to\bbbar$ dominates
the inclusive $\bbbar$ rate. Our example switches it off
(\code{HardQCD:gg2bbbar = off}) to keep a single channel whose matrix element
fits on one line. Nothing in the rest of the journey depends on this choice,
but a reader comparing with an inclusive $\bbbar$ sample must remember it.

The sum over $q$ in \cref{eq:ancestral_home:factorisation} runs over the
flavours \PYTHIA allows in the initial state, \code{PDFinProcess:nQuarkIn=5}
by default: $d,u,s,c$ and $b$. For $q=b$ the process $b\bar b\to b\bar b$ has
additional $t$-channel diagrams that \cref{eq:ancestral_home:dsigma} ignores;
\PYTHIA process~124 applies the $s$-channel formula to all five flavours, and
so do our integrators. This is a convention we reproduce, not a complete
calculation. It matters at the $2\%$ level, the size of the incoming-$b$
term (\cref{tab:ancestral_home:flavours}).

Treating the $b$ as a massless parton inside the proton while giving it a
mass in the final state is the \emph{five-flavour scheme} (5FS). The
alternative \emph{four-flavour scheme} (4FS) keeps the $b$ out of the proton
and generates it only in the matrix element, at the price of large logarithms
$\ln(\pthat^2/\mb^2)$ that the 5FS resums through PDF evolution. For $\bbbar$
production at $\pthat\ge30\GeV$ the two schemes must agree once higher orders
are included; at leading order they do not, and the 5FS with
\cref{eq:ancestral_home:dsigma} is simply what \PYTHIA and CMS use.

\subsection{Scales and the coupling}
\label{sec:ancestral_home:scales_and_coupling}

Every leading-order calculation has to choose $\muR$ and $\muF$. \PYTHIA's
default for a $2\to2$ process with two equal-mass final particles is the
transverse mass,
\begin{equation}
  \muR^2 = \muF^2 = \pthat^2 + \mb^2 ,
  \label{eq:ancestral_home:scale_choice}
\end{equation}
which its manual calls scale codes~1--3 (they coincide for equal masses);
code~4 uses $\shat$ and code~5 a fixed value. The multipliers
\code{SigmaProcess:renormMultFac} and \code{factorMultFac} multiply the
\emph{squared} scale, a source of factor-of-two confusion when one tries to
``double the scale''.

The coupling in the hard process is a separate object from the coupling in
the parton shower and from the one stored with the PDF set. The CP5
tune~\cite{CMS:2019csb} sets \code{SigmaProcess:alphaSvalue = 0.118} with
second-order running (\code{alphaSorder = 2}) for the matrix element, and the
same for the initial- and final-state showers and for multi-parton
interactions. Using an NNLO PDF set does not make the matrix element NNLO,
and using the PDF set's own $\alphas$ would not reproduce \PYTHIA. At our
event the running gives $\alphas(1178\GeV^2)=0.1385$, well above
$\alphas(m_Z^2)=0.118$ because the scale is below $m_Z$.

\section{From formula to number: the cross section step by step}
\label{sec:ancestral_home:from_formula_to_number}

We now have everything needed to turn \cref{eq:ancestral_home:factorisation}
into $18.5\nb$. The steps below are the steps of the standalone code in
\cref{sec:ancestral_home:standalone_code}; the function names in brackets
point to it.
\begin{enumerate}
  \item \textbf{Choose the integration variables.} Map the unit cube
    $(u_1,u_2,u_3)\in[0,1]^3$ to $\ln\tau$, $Y$ and $|\cos\theta|$:
    \begin{align*}
      \ln\tau &= \ln\tau_{\min} + u_1\,\ln(1/\tau_{\min}), &
      Y &= (2u_2-1)\,Y_{\max}, \quad Y_{\max}=-\tfrac12\ln\tau, \\
      |\cos\theta| &= u_3\,c_{\max}, &
      c_{\max} &= \sqrt{1-(\pthat^{\min}/p^{*})^2}.
    \end{align*}
    Only positive $\cos\theta$ is sampled, with a factor~2, because the
    matrix element and all our cuts are even in $\cos\theta$
    [\code{born\_weights}].
  \item \textbf{Compute the kinematics.} $\shat=\tau s$, $x_{1,2}$ from
    \cref{eq:ancestral_home:tauY}, $\pthat^2$ from
    \cref{eq:ancestral_home:kinematics}, and the scale
    $Q^2=\pthat^2+\mb^2$.
  \item \textbf{Evaluate the ingredients at that point.} $\alphas(Q^2)$
    [\code{TableAlphaS}], the ten PDF values $x_1f_q$, $x_1f_{\bar q}$,
    $x_2f_q$, $x_2f_{\bar q}$ for $q=d,u,s,c,b$ [\code{TablePDF.xf}], and
    the matrix element $\mathrm{d}\hat\sigma/\mathrm{d}\cos\theta$
    [\code{dsigma\_dcostheta}].
  \item \textbf{Form the weight.} The integrand times the volume of the mapped
    region,
    \[
      w_q = \ln\!\frac{1}{\tau_{\min}}\;\cdot\;2Y_{\max}\;\cdot\;2c_{\max}\;\cdot\;
      \frac{\mathrm{d}\hat\sigma}{\mathrm{d}\cos\theta}\;
      \bigl[\,x_1f_q(x_1)\,x_2f_{\bar q}(x_2)+x_1f_{\bar q}(x_1)\,x_2f_q(x_2)\,\bigr],
    \]
    one term per incoming flavour, both beam orientations included. The
    cross section is the average of $\sum_q w_q$ over the cube.
  \item \textbf{Average over many points.} Quasi-random Sobol points
    ($2^{16}=65\,536$ per pass) fill the cube far more evenly than
    pseudo-random ones, so a single pass already gives four significant digits
    [\code{integrate}].
  \item \textbf{Estimate the numerical error.} Repeat with independently
    scrambled Sobol sequences and take the standard deviation of the pass
    totals divided by $\sqrt{N_{\mathrm{repeats}}}$. The flavour terms are
    correlated, so the total is summed before the spread is taken.
  \item \textbf{Convert units.} Multiply by $3.893\,793\,72\times10^{8}$ to
    go from $\GeV^{-2}$ to pb.
\end{enumerate}

The results are in \cref{tab:ancestral_home:three_numbers} and, per flavour, in
\cref{tab:ancestral_home:flavours}. Up quarks dominate because $u$ is the
most abundant valence quark and the $\bar u$ sea is comparable to $\bar d$.
The dependence on the \pt cut (\cref{tab:ancestral_home:ptscan}) is steep:
raising $\pthat^{\min}$ from $30$ to $100\GeV$ costs a factor~54.

\begin{table}[htbp]\centering
\caption[Contribution of each incoming flavour]{Contribution of each incoming flavour to the default cross section
(eight repeats of $2^{16}$ Sobol points, \file{standalone_repeats8.out}).}
\label{tab:ancestral_home:flavours}
\begin{tabular}{lrr}\toprule
Incoming pair & $\sigma$ [nb] & fraction \\ \midrule
$u\bar u$ & $8.4089$ & $45.4\%$ \\
$d\bar d$ & $6.1746$ & $33.4\%$ \\
$s\bar s$ & $2.4380$ & $13.2\%$ \\
$c\bar c$ & $1.0818$ & $5.8\%$ \\
$b\bar b$ & $0.4032$ & $2.2\%$ \\ \midrule
total & $18.5065 \pm 0.00005$ & \\
\bottomrule
\end{tabular}
\end{table}

\begin{table}[htbp]\centering
\caption[Dependence on the generator-level cut $\pthat^{\min}$]{Dependence on the generator-level cut $\pthat^{\min}$, all other
settings at their defaults (\file{example_pt_scan.json}). The \PYTHIA
reference of $18.44\nb$ applies to the $30\GeV$ row only.}
\label{tab:ancestral_home:ptscan}
\begin{tabular}{rr}\toprule
$\pthat^{\min}$ [GeV] & $\sigma$ [nb] \\ \midrule
20 & $62.086$ \\
30 & $18.507$ \\
50 & $3.6505$ \\
100 & $0.34532$ \\
\bottomrule
\end{tabular}
\end{table}

\subsection{What this number is, and what it is not}
\label{sec:ancestral_home:caveats}

The agreement in \cref{tab:ancestral_home:three_numbers} was obtained without
fitting or adjusting anything to the \PYTHIA result. It holds because the
integrator reproduces \PYTHIA's conventions, and each convention is a caveat
the reader should carry into the later chapters:
\begin{itemize}
  \item It is the \emph{Born} hard-process rate, unweighted and unsuppressed.
    Nothing has showered, hadronised or been clustered. It is not a $b$-jet
    cross section, and a cut on $\pthat$ is not a cut on jet \pt.
  \item The $|y|$ or $|\eta|$ cuts the integrator offers act on the two Born
    quarks, not on reconstructed jets.
  \item The incoming-$b$ term uses the $s$-channel formula only
    (\cref{sec:ancestral_home:gg_and_flavours}); the $gg$ channel is absent by
    choice.
  \item PDFs are frozen at grid edges, clipped at zero, and the heavy sea is
    symmetrised (\cref{sec:ancestral_home:pdfs}).
  \item The coupling is \PYTHIA's second-order $\alphas$ with
    $\alphas(m_Z)=0.118$, not the PDF set's.
  \item The quoted errors are integration or sampling errors. They say nothing
    about the accuracy of leading-order QCD, which is the subject of
    \cref{sec:ancestral_home:scales}.
  \item Points in a \pt scan share the same Sobol seeds and are therefore not
    statistically independent bins.
  \item User hooks, process weights, phase-space biasing, nuclear PDFs and
    any post-shower selection that a production configuration might add are
    not reproduced.
\end{itemize}

\section{Standalone code: the whole calculation in one file}
\label{sec:ancestral_home:standalone_code}

The file \file{example/ME/Standalone/standalone_qqbar_bbbar_xsec.py} contains
the calculation of \cref{sec:ancestral_home:from_formula_to_number} and
nothing else. It needs Python~3.10 or newer with NumPy and SciPy; no LHAPDF,
no \PYTHIA. The PDF grid (member~0 of the NNPDF3.1 set on LHAPDF's own knots,
restricted to the phase space of the default run) and \PYTHIA's $\alphas$
curve are embedded as numeric tables between the markers
\code{BEGIN/END GENERATED TABLES}, produced by
\file{dump_standalone_tables.py} in an environment that has both libraries.
From the repository root:
\begin{minted}[linenos=false,frame=single,framesep=1.5mm,fontsize=\small]{bash}
python3 example/ME/Standalone/standalone_qqbar_bbbar_xsec.py              # one pass
python3 example/ME/Standalone/standalone_qqbar_bbbar_xsec.py --repeats 8  # with error
python3 example/ME/Standalone/standalone_qqbar_bbbar_xsec.py --ptmin 50 --ptmax 100
\end{minted}
The first command takes about a second and prints the output reproduced
below. The two integrators agree to $5\times10^{-6}$ for the same Sobol seed,
so the tables lose nothing; if the PDF set, the coupling or the kinematic range
is changed the tables must be regenerated, and the script says so.

\filelisting{text}{example/ME/Standalone/standalone_default.out}{output of the first command above}

\subsection{Reading the code}

The listing follows. A map for the reader: \code{TablePDF} (bicubic spline
in $(\ln x,\ln Q^2)$, with the freezing, clipping and heavy-sea conventions of
\cref{sec:ancestral_home:pdfs}); \code{TableAlphaS} (monotone cubic in
$\ln Q^2$, refusing to extrapolate); \code{dsigma\_dcostheta}
(\cref{eq:ancestral_home:dsigma}, vectorised); \code{born\_weights} (steps~1
to~4 of \cref{sec:ancestral_home:from_formula_to_number}, including the
numerically stable form of $c_{\max}-c_{\min}$ when an upper \pt cut is
present); \code{integrate} (steps~5 and~6); and \code{main} (command line,
sanity checks and the printout). The line holding the tables is replaced by a
placeholder in print; the file in the repository has the numbers.

\IfFileExists{build/standalone_qqbar_bbbar_xsec.py}{%
  \filelisting{python}{build/standalone_qqbar_bbbar_xsec.py}{print copy; the table line is a placeholder}}{%
  \todo{Run \code{make listings} to generate the print copy of the standalone script.}}

\subsection{The same calculation with LHAPDF and \PYTHIA}

The sibling script \file{qqbar_bbbar_xsec.py} in the same directory does the
same integral with LHAPDF for the PDFs and \PYTHIA's own \code{AlphaStrong} class
for the coupling, and offers every knob discussed in this chapter as a command
line option: upper and lower $\pthat$ and $\hat m$ cuts, $|y|$ or $|\eta|$
cuts on both quarks, the number of incoming flavours, all five \PYTHIA scale
codes and their multipliers, a $K$ factor, the PDF member, the coupling
backend and order, and the option to keep negative PDF values. Its
\file{README.md} explains how to build the environment and how to export the
actual settings of a running \PYTHIA instance with the header
\file{export_pythia_xsec.h}, so that one integrates precisely what the
generator was configured to do rather than what a log file suggests. We use it
in \cref{sec:ancestral_home:scales} for the variations that the standalone
script does not expose.

\section{How CMS does it in practice: the GEN step in \CMSSW}
\label{sec:ancestral_home:cmssw}

In CMS the hard process is not run by hand; it is one step of a production
chain driven by \code{cmsDriver.py}. The configuration in
\file{example/ME/CMSSW/step1_GEN_cfg.py} was generated with the command in
its header, here broken into lines:
\begin{minted}[linenos=false,frame=single,framesep=1.5mm,fontsize=\small,breaklines]{bash}
cmsDriver.py Configuration/GenProduction/python/BBbar_14TeV_TuneCP5_cfi.py \
  --conditions auto:phase1_2024_realistic --era Run3_2024 --geometry DB:Extended \
  -n 1 -s GEN --datatier GEN --eventcontent FEVTDEBUG \
  --fileout file:output_step1_GEN.root --python_filename step1_GEN_cfg.py \
  --no_exec --beamspot DBrealistic \
  --customise_commands "process.generator.comEnergy = cms.double(13600.0)"
\end{minted}
The pieces are: a \emph{fragment} (the \code{\_cfi.py} file, which carries the
\PYTHIA settings), the detector \emph{conditions} and \emph{era}, the
\emph{step} (\code{GEN} only), the output data tier and event content, and a
customisation that overrides the fragment's $14\TeV$ with the Run~3 energy
of $13.6\TeV$. To run it:
\begin{minted}[linenos=false,frame=single,framesep=1.5mm,fontsize=\small]{bash}
cmsrel CMSSW_15_0_5
cd CMSSW_15_0_5/src
cmsenv
cmsRun step1_GEN_cfg.py
\end{minted}

\subsection{Reading the configuration}

Most of the file is boilerplate that every \CMSSW job carries. The physics
is in the module \code{process.generator}, an \code{EDFilter} of type
\code{Pythia8ConcurrentGeneratorFilter}, and in particular in three blocks
of its \code{PythiaParameters}:
\begin{itemize}[nosep]
  \item \code{processParameters}: the three lines that define our process,
    \code{HardQCD:gg2bbbar = off}, \code{HardQCD:qqbar2bbbar = on} and
    \code{PhaseSpace:pTHatMin = 30.}. Everything in
    \cref{sec:ancestral_home:matrix_element} follows from these.
  \item \code{pythia8CP5Settings}: the CMS tune~\cite{CMS:2019csb}. It fixes
    the PDF through \code{PDF:pSet} to the NNPDF3.1 set of
    \cref{sec:ancestral_home:pdfs}, and sets the four couplings, for the matrix
    element, the two showers and the multi-parton interactions, to
    $\alphas(m_Z)=0.118$ at second order. It also carries the
    multi-parton-interaction and colour-reconnection parameters that later
    chapters will meet.
  \item \code{pythia8CommonSettings}: production defaults such as the
    tolerance on energy--momentum conservation and the lifetime cut-off for
    particle decays.
\end{itemize}
Other lines worth noticing: \code{maxEvents.input = 1000} (the fragment
name says one event, the job was run with a thousand); the
\code{comEnergy} override at the end; the output module writing
\code{FEVTDEBUG} content to \file{output_step1_GEN.root}; and a commented-out
\code{customiseJitJet} line that hooks in the book's own analysis package,
used from the next chapter on. The full file:

\filelisting{python}{example/ME/CMSSW/step1_GEN_cfg.py}{\code{cmsDriver.py} configuration for the GEN step}

\subsection{Reading the log}

\PYTHIA writes two things to \file{step1_GEN_cfg.log} that close the loop
with \cref{sec:ancestral_home:from_formula_to_number}. For the first event it
prints the ``Info Listing'' with the incoming partons, their $x$ and PDF
values, the partonic invariants and the coupling; these are the numbers of
\cref{tab:ancestral_home:pdf_values} and \cref{eq:ancestral_home:kinematics}.

\filelistinglines[\widelisting]{text}{example/ME/CMSSW/step1_GEN_cfg.log}{187}{197}{the hard-process kinematics of our event}

At the end of the job it prints the cross section it estimated while
sampling. Three thousand phase-space points were tried to accept one thousand
unweighted events; the cross section is the average weight of the tried
points, not the acceptance ratio.

\filelistinglines[\widelisting]{text}{example/ME/CMSSW/step1_GEN_cfg.log}{237}{248}{cross-section statistics; PYTHIA reports in mb}

The value $1.833\times10^{-5}\,\mathrm{mb}=18.33\nb$ with a statistical
error of $0.30\nb$ is $0.6$ standard deviations from our $18.507\nb$; a
longer standalone \PYTHIA~8.315 run gave $18.44\pm0.10\nb$, $0.7$ standard
deviations away. The residual $0.4\%$ is consistent with sampling
fluctuations and with the \PYTHIA version difference (8.311 in \CMSSW versus
8.315 for the coupling table), and nothing in it calls for a correction.

\subsection{The same calculation in \MG}
\label{sec:ancestral_home:madgraph}

A reader who generates Event~A with \MG~\cite{Alwall:2014hca} at leading order
computes the same Born matrix element, plus the $gg$ channel unless it is
excluded. The correspondence is: \code{generate p p > b b\~{}} in the process
card (with \code{define p = u d s c b u\~{} d\~{} s\~{} c\~{} b\~{}} for the
five-flavour proton and the model \code{sm-no\_b\_mass} or \code{loop\_sm}
variants depending on whether the $b$ is massive); in the run card
\code{ebeam1 = ebeam2 = 6800}, \code{pdlabel = lhapdf}, \code{lhaid = 303600},
\code{ptb = 30} for the $b$ \pt cut, and \code{dynamical\_scale\_choice}
together with \code{scalefact} for the scale. \MG stores the result as
weighted or unweighted events in the Les Houches Event (LHE) format, whose
\code{<event>} blocks carry the same information as \PYTHIA's hard-process
listing in \cref{sec:ancestral_home:our_jet}; the plotting script of that
section reads both formats.

\section{Our jet at this stage: the \texorpdfstring{$b$}{b}-quark four-vector in the hard-process record}
\label{sec:ancestral_home:our_jet}

Here is the mother at the moment of her birth. The file
\file{example/ME/event_ME.lhe} holds \PYTHIA's hard-process listing for the
event followed through the book, copied from the log of the \CMSSW job:

\filelisting[\widelisting]{text}{example/ME/event_ME.lhe}{the hard-process record of our event; momenta and masses in GeV}

Reading the columns: \code{status} $-12$ marks the two beam protons, $-21$ the
two incoming partons after they have been resolved out of the protons, and
$23$ the two outgoing partons of the hard process; negative status means the
particle has been ``used up'' and positive means it is still there for the
next stage. The colour columns show the colour flow: the $\bar u$ carries
anticolour~101, which is handed to the $\bar b$, and the $u$ carries
colour~102, handed to the $b$; the gluon in between never appears because it is
internal. The momentum sum at the bottom is that of the two incoming partons,
$p_z=1.871-1931.557=-1929.69\GeV$: the whole system recoils backward.

\begin{table}[htbp]\centering
\caption[Kinematics of the two quarks of our event]{Kinematics of the two quarks in \file{event_ME.lhe}, computed by
\file{example/ME/plot_journey_ME.py}. The $b$ quark is the mother we follow.}
\label{tab:ancestral_home:kinematics}
\begin{tabular}{lrrrrrrrr}\toprule
 & $p_x$ & $p_y$ & $p_z$ & $E$ & $\pt$ & $\eta$ & $y$ & $\phi$ \\
 & [GeV] & [GeV] & [GeV] & [GeV] & [GeV] & & & [rad] \\ \midrule
$b$       & $-18.630$ & $-28.420$ & $-171.125$  & $174.533$  & $33.98$ & $-2.32$ & $-2.31$ & $-2.151$ \\
$\bar b$  & $18.630$  & $28.420$  & $-1758.560$ & $1758.895$ & $33.98$ & $-4.64$ & $-4.63$ & $0.991$ \\ \midrule
pair      & $0$ & $0$ & $-1929.685$ & $1933.428$ & $0$ & & $Y=-3.47$ & \\
\bottomrule
\end{tabular}
\end{table}

\Cref{tab:ancestral_home:kinematics} and \cref{fig:ancestral_home:journey}
show what the record means. The two quarks are exactly back to back in the
transverse plane ($\Delta\phi=\pi$) with $\pt=33.98\GeV$ each, just above
the $30\GeV$ generator cut, as most events in a steeply falling spectrum are.
The pair mass is $m_{\bbbar}=\sqrt{\shat}=120.2\GeV$, and
$\shat=x_1x_2s=2.752\times10^{-4}\times0.2841\times13600^2=14\,456\GeV^2$
checks against the log. The pair rapidity $Y=-3.47$ and the decay
rapidity $y^{*}=\operatorname{artanh}(\beta\cos\hat\theta)=-1.16$ give
$y_b=Y-y^{*}=-2.31$ and $y_{\bar b}=Y+y^{*}=-4.63$, in agreement with the
table. The mother, the $b$, is born at $\eta=-2.32$: inside the tracker
acceptance of $|\eta|<2.5$, but only just. The $\bar b$ at $\eta=-4.64$ heads
into the forward hadron calorimeter, where there are no tracks and therefore,
much later in the journey, no $b$ tag. Two quarks of identical \pt will have
completely different fates.

\begin{figure}[htbp]\centering
\includegraphics[width=\textwidth]{ME/journey_ME.pdf}
\caption[The hard process of our event in the $(\eta,\phi)$ plane]{The hard process of our event in the $(\eta,\phi)$ plane, drawn by
\file{example/ME/plot_journey_ME.py} from \file{example/ME/event_ME.lhe}.
Marker area is proportional to \pt; the shaded bands are the CMS tracker
region, the endcap calorimeters and the forward calorimeter (HF). The incoming
partons and their momentum fractions are noted at the two edges. Nothing has
showered, hadronised or reached a detector yet.}
\label{fig:ancestral_home:journey}
\end{figure}

The ledger opens with a single entry: one parton, its \pt, its mass. Every
later row will be compared with this one.

\begin{ledger}{The Ancestral Home}
  \ledgerrow{before this stage}{--}{--}{--}{--}
  \ledgerrow{$b$ quark in the hard-process record}{33.98}{4.80}{1}{--}
\end{ledger}

\section{Scale choices and their uncertainties}
\label{sec:ancestral_home:scales}

A leading-order cross section depends on $\muR$ through $\alphas^2$ and on
$\muF$ through the PDFs, and neither dependence is compensated by anything at
this order. The conventional estimate of the missing higher orders is to
vary both scales by factors of two around the central choice
\cref{eq:ancestral_home:scale_choice}. \Cref{tab:ancestral_home:variations}
shows the result for our process, computed with the full integrator.

\begin{table}[htbp]\centering
\caption[Scale, flavour and PDF-convention variations]{Variations of the default calculation, computed with
\file{qqbar_bbbar_xsec.py} ($2^{15}$ Sobol points, two repeats; integration
errors below $0.002\%$). The central value is $18.507\nb$.}
\label{tab:ancestral_home:variations}
\begin{tabular}{llrr}\toprule
Variation & Option & $\sigma$ [nb] & change \\ \midrule
$\muR\times2$              & \code{--mur-factor 2}                   & $14.765$ & $-20.2\%$ \\
$\muR\times\tfrac12$       & \code{--mur-factor 0.5}                 & $23.941$ & $+29.4\%$ \\
$\muF\times2$              & \code{--muf-factor 2}                   & $21.204$ & $+14.6\%$ \\
$\muF\times\tfrac12$       & \code{--muf-factor 0.5}                 & $15.621$ & $-15.6\%$ \\
$\muR,\muF\times2$         & \code{--mur-factor 2 --muf-factor 2}     & $16.911$ & $-8.6\%$ \\
$\muR,\muF\times\tfrac12$  & \code{--mur-factor 0.5 --muf-factor 0.5} & $20.198$ & $+9.1\%$ \\ \midrule
no incoming $b$            & \code{--nquark 4}                       & $18.104$ & $-2.2\%$ \\
negative PDFs kept         & \code{--signed-pdfs}                    & $18.498$ & $-0.05\%$ \\
\bottomrule
\end{tabular}
\end{table}

Three lessons are in the table. First, the scale uncertainty of this
leading-order number is of order $20$ to $30\%$, two orders of magnitude
larger than any statistical error in
\cref{tab:ancestral_home:three_numbers}; agreement between two leading-order
codes at $0.4\%$ says nothing about agreement with nature. Second, $\muR$ and
$\muF$ pull in opposite directions: a larger $\muR$ lowers $\alphas$, while a
larger $\muF$ raises the small-$x$ sea that our $\bar u$ came from, because
DGLAP evolution feeds the sea from the gluon as $Q^2$ grows. Varying both
together therefore under-estimates the uncertainty, which is why the standard
prescription is the envelope of a seven-point variation
($\muR,\muF$ each $\times\tfrac12,1,2$ without the two extreme opposite
combinations). Third, the two ``convention'' rows are small: whether negative
PDF values are clipped is irrelevant, and the incoming-$b$ approximation of
\cref{sec:ancestral_home:gg_and_flavours} is a $2\%$ effect.

Beyond scales, a complete uncertainty would add the PDF uncertainty from the
other members of the set (100 replicas for NNPDF3.1, run with
\code{--member}), the variation of $\alphas(m_Z)$ by $\pm0.001$, and the
choice of $\mb$, which enters through $\rho$ and through the scale. The real
cure is the next order: at NLO the $\muR$ and $\muF$ dependence of the Born
term is cancelled by the explicit logarithms of the one-loop and real-emission
corrections, and the residual uncertainty shrinks to about $10\%$. That is what
\MG at NLO~\cite{Alwall:2014hca} provides, and why CMS uses NLO samples where
the cross section itself matters. For following one jet through the detector,
leading order is enough; for a measurement it is not.

\section{Further reading}
\label{sec:ancestral_home:further_reading}
The factorisation theorem and the massive-quark Born formula are reviewed in
the Particle Data Group's ``Cross-section formulae'' and ``QCD'' sections
\cite{ParticleDataGroup:2024cfk}. NNPDF3.1 is described in
\cite{NNPDF:2017mvq}, the LHAPDF library in \cite{Buckley:2014ana}, and the
PDF4LHC recommendations for combining PDF uncertainties in
\cite{Butterworth:2015oua}. The \PYTHIA~8.3 guide~\cite{Bierlich:2022pfr},
chapters on hard processes and on couplings and scales, is the reference for
process~124 and the scale codes; the CP5 tune is documented in
\cite{CMS:2019csb}. \MG and its NLO capabilities are described in
\cite{Alwall:2014hca}.

\begin{pitfalls}
\begin{itemize}[nosep]
  \item \textbf{Three numbers, one cross section.} $18.51$, $18.44$ and
    $18.33\nb$ are the same calculation with different statistical
    fluctuations. Do not ``fix'' one to match another, and do not read the
    $0.4\%$ spread as physics: the theoretical uncertainty is $20$--$30\%$.
  \item \textbf{$\pthat$ is not jet \pt.} The generator cut acts on the two
    Born quarks. After showering, a $34\GeV$ quark becomes a jet of any \pt
    from a few GeV up; a sample with $\pthat\ge30\GeV$ is not a sample of
    jets above $30\GeV$, and cannot be used for jets near that threshold.
  \item \textbf{NNLO PDF, LO matrix element.} The order of the PDF set and the
    order of the calculation are independent choices. CMS uses NNLO PDFs with
    LO \PYTHIA processes; that is a tune decision, not a claim of NNLO accuracy.
  \item \textbf{Scale multipliers multiply $Q^2$.} \PYTHIA's
    \code{renormMultFac = 4} doubles $\muR$, not quadruples it. The
    integrator's \code{--mur-factor} option multiplies $\muR$ itself.
  \item \textbf{Which $\alphas$?} The matrix element, each shower, the
    multi-parton interactions and the PDF set each have their own coupling.
    Check \code{SigmaProcess:alphaSvalue} and \code{alphaSorder}, not
    \code{Tune:pp}.
  \item \textbf{$\mb=4.8\GeV$} is \PYTHIA's default mass for the $b$ in the
    matrix element. It is neither the pole mass nor the $\overline{\mathrm{MS}}$
    mass in the Particle Data Group tables, and changing it changes the
    threshold, the scale and the cross section.
  \item \textbf{$gg$ is off here.} An inclusive $\bbbar$ sample is
    dominated by gluon fusion. Our number covers only $q\bar q$ annihilation.
  \item \textbf{The \file{.lhe} file is not an LHE file.} It is \PYTHIA's
    hard-process listing, saved with that extension because it plays the role
    of the LHE record. Status codes ($-21$, $23$) are \PYTHIA's, not the LHE
    standard's ($-1$, $1$).
  \item \textbf{Units.} \PYTHIA prints mb, the integrators print pb and nb.
    $1\,\mathrm{mb}=10^{6}\nb=10^{9}\pb$.
  \item \textbf{Which Python?} A shell alias such as
    \code{alias python=/usr/bin/python3} silently overrides a virtual
    environment. Run the scripts with the explicit interpreter path when in
    doubt; the standalone script prints the table provenance on its first
    line so the reader can see what was loaded.
\end{itemize}
\end{pitfalls}

\begin{exercises}
\item \textbf{Run it.} Run \file{standalone_qqbar_bbbar_xsec.py} with no
  options and confirm $18.5064\nb$. Run it with \code{--repeats 8} and
  compare the error with \cref{tab:ancestral_home:three_numbers}. The
  \CMSSW job accepted $1000$ unweighted events and quoted a $1.7\%$
  statistical error; how many events would reach $0.1\%$? Why is that
  precision pointless for a leading-order calculation?
\item \textbf{The angular integral by hand.} Integrate
  \cref{eq:ancestral_home:dsigma} over $\cos\theta$ from $-1$ to $1$ and
  obtain \cref{eq:ancestral_home:sigmahat}. Evaluate both at our event
  ($\shat=14\,456\GeV^2$, $\alphas=0.1385$, $\cos\hat\theta=-0.824$) and
  confirm $\hat\sigma=481\pb$ and
  $\mathrm{d}\hat\sigma/\mathrm{d}\cos\theta=302\pb$. Then integrate only over
  $|\cos\theta|\le c_{\max}$ for $\pthat^{\min}=30\GeV$: what fraction of
  $\hat\sigma$ survives the cut at this $\shat$?
\item \textbf{Kinematics from the record.} Using only the four-vectors in
  \file{event_ME.lhe} and the beam energy, compute $x_1$, $x_2$, $\shat$, $Y$,
  $\pthat$, $\cos\hat\theta$, $y^{*}$ and the rapidities and pseudorapidities
  of both quarks. Compare with \cref{tab:ancestral_home:kinematics} and with
  the printout of \file{plot_journey_ME.py}. Why do $y$ and $\eta$ differ by
  $0.01$ here, and when would they differ by much more?
\item \textbf{The \pt spectrum.} Reproduce \cref{tab:ancestral_home:ptscan}
  with \code{--ptmin 20}, \code{50} and \code{100}. Between $50$ and
  $100\GeV$ the cross section falls as $(\pthat^{\min})^{-n}$; determine $n$.
  Dimensional analysis of \cref{eq:ancestral_home:sigmahat} suggests
  $n\approx4$; explain what the PDFs and the running coupling do to it.
\item \textbf{Scale variation with one line.} In \code{born\_weights} the
  line \code{q2 = pt2 + mb2} sets both scales. Change it to
  \code{4*(pt2 + mb2)} and to \code{0.25*(pt2 + mb2)} and reproduce the two
  ``both scales'' rows of \cref{tab:ancestral_home:variations}. Why can the
  standalone script not reproduce the single-scale rows without a larger
  change? (Look at what \code{TablePDF.xf} and \code{TableAlphaS} are called
  with.)
\item \textbf{Where does the $\muF$ dependence come from?} Using the
  embedded PDF table, evaluate $x f_{\bar u}(x,Q^2)$ at $x=2.75\times10^{-4}$
  and $x f_u$ at $x=0.284$ for $Q^2=1178/4$, $1178$ and $4\times1178\GeV^2$.
  Which of the two moves more, and in which direction? Connect this to the
  sign of the $\muF$ rows in \cref{tab:ancestral_home:variations}.
\item \textbf{The incoming $b$.} Zero the fifth flavour in
  \code{born\_weights} (or run the full script with \code{--nquark 4}) and
  confirm the $2.2\%$ change. Draw the additional Feynman diagrams that
  $b\bar b\to b\bar b$ has and $u\bar u\to b\bar b$ does not. Why does the
  same-flavour case not matter much here, and when would it?
\item \textbf{Conventions.} The default full-environment run clipped
  $1\,166\,234$ negative PDF values yet the cross section moved by $0.05\%$
  with \code{--signed-pdfs}. Where in the $(x,Q^2)$ plane and for which
  flavours do fitted PDFs go negative, and why does it not matter for this
  process?
\item \textbf{The gluon channel (needs \CMSSW).} In \file{step1_GEN_cfg.py}
  switch the gluon channel on, \code{HardQCD:gg2bbbar = on}, run $1000$ events
  and read the per-process cross sections from the statistics table in the
  log. What fraction of the inclusive $\bbbar$ rate at $\pthat\ge30\GeV$ is
  $q\bar q$ annihilation? Why is the $gg$ matrix element not a one-line
  formula?
\item \textbf{Two fates.} The $\bar b$ of our event is at $\eta=-4.64$.
  Which CMS subdetectors will see it, and which will not? Will it become a
  jet in \cref{ch:the_inn}, and can it ever be $b$-tagged? Keep your answer;
  \cref{ch:registry_office} will check it.
\end{exercises}
